\chapter{Conclusions and Outlook} \label{ch:conclusion}

This final chapter evaluates the overarching implications of the research conducted in this thesis. Rather than summarizing the quantitative results, it extracts the core lessons learned and distills the complex strategic interactions into a clear, fundamental takeaway.

The primary methodological achievement of this thesis is the successful development and stabilization of an Equilibrium Problem with Equilibrium Constraints (EPEC) tailored to Battery Energy Storage Systems (BESS) allocation. The results clearly demonstrate that the framework works: It successfully shifts the modeling paradigm from a naive central-planner cost-minimization approach to a decentralized, game-theoretic oligopoly. By reaching a stable Nash Equilibrium, the model proves that capturing the strategic self-interest of independent, profit-maximizing investors is both computationally feasible and absolutely essential for realistic energy system modeling.

However, the application of this framework revealed several unexpected realities and necessary compromises that fundamentally challenge traditional BESS allocation modeling:
\begin{itemize}
    \item \textbf{Aggressive Siting:} Investors are driven by the profitability of system-wide markets and therefore aggressively distribute their assets across the entire grid to exhaust all available nodal connection limits and maximize their individual market shares.
    \item \textbf{Arbitrage-Reserve Trade-off:} The most significant operational finding is that lucrative capacity premiums in the aFRR market completely override traditional day-ahead logic. Investors lock up their expensive inverter capacity (MW) for reserve provision, creating a severe power bottleneck. To still participate in the spot market, they massively oversize their cheap energy capacity (MWh).
    \item \textbf{Methodological Pragmatism:} The continuous, non-linear nature of EPECs required significant engineering pragmatism. To prevent interior-point solvers from crashing at mathematical singularities or falling into oscillatory traps, relaxing KKT complementarity conditions and implementing a residual demand logic were mandatory compromises.
    \item \textbf{Future Evolution:} To overcome these continuous approximations, future research must aim towards Mixed-Integer Programs with Equilibrium Constraints (MIPEC). This would allow for the integration of discrete bidding steps, true Pay-as-Bid clearing, and non-differentiable battery degradation algorithms (e.g., Rainflow-Counting).
\end{itemize}

Ultimately, the most profound insight of this thesis extends beyond the mathematical equations. If one were to summarize the findings of this research in a single takeaway, it is this: \textit{Optimal battery sizing and operation are not dictated by the physical availability of renewable energy, but by the regulatory design of the markets they participate in.} 

When ancillary service markets are highly lucrative, strategic investors will act as rational opportunists. They will willingly accept inefficient, slow-cycling, and even loss-making operations on the spot market, simply because securing a dominant market share in the reserve market compensates for these opportunity costs. Therefore, evaluating battery investments based on single-market business cases or isolated grid metrics is an insufficient approach. To understand the true dimensioning and operational logic of storage systems in modern power grids, it is absolutely essential to account for the severe cross-market spillover effects driven by strategic investor behavior.





% This final chapter evaluates the overarching implications of the research conducted in this thesis. Rather than providing a mere summary of the quantitative results presented in Chapter \ref{ch:results}, it critically assesses the validity of the research question, the appropriateness of the chosen methodology, and the core lessons learned. Finally, it outlines how the identified limitations can be overcome and provides a comprehensive outlook for future research in the field of energy system modeling.

% As introduced in Chapter \ref{ch:introduction}, the guiding research question of this thesis asked what insights could be gained by analyzing BESS allocation through a game-theoretic EPEC framework rather than a traditional central-planner approach. The transition towards weather-dependent renewable energy sources necessitates massive storage capacities. However, the high capital intensity of BESS requires precise, market-driven investment strategies rather than simplified, purely grid-oriented heuristics. The results of this thesis undeniably confirm the relevance of this research question: The organic distribution of BESS and their respective Energy-to-Power ratios are not driven by simple renewable availability or systemic cost-minimization, but heavily depend on complex market interactions, capital costs, and nodal grid limitations. Therefore, shifting the methodological paradigm to capture these strategic realities was absolutely essential.

% To answer this research question, the methodology of an Equilibrium Problem with Equilibrium Constraints was chosen and mathematically formulated in Chapter \ref{ch:Methodology}. While traditional central-planner optimization models (such as pure LOPF) provide valuable systemic benchmarks, they entirely fail to capture the strategic self-interest and market power of independent actors in liberalized electricity markets. Despite the immense mathematical challenges and necessary algorithmic compromises -- such as KKT relaxation and the implementation of a residual demand framework -- the EPEC approach proved to be highly appropriate. It was the only methodology capable of realistically depicting oligopolistic competition, endogenous price formation via continuous Cournot Inverse Demand curves, and the subsequent strategic behavior of investors competing for shared nodal capacities.

% The synthesis of the results in Chapter \ref{ch:synthesis} yielded several fundamental insights into strategic BESS investments. Regarding the \textit{siting}, the framework demonstrated that investors organically distribute their capacities across the network to avoid nodal congestion, thereby actively stabilizing the grid without the need for central-planner redispatch interventions. Regarding the \textit{sizing}, the optimal MWh capacity was revealed to be a direct function of the investors' WACC and the strict operational requirements of the reserve market. The most significant finding, however, is a severe Arbitrage-Reserve Trade-off. The co-optimization of day-ahead spot and aFRR reserve markets proved that highly lucrative capacity premiums in the reserve market can entirely override traditional spot market logic. This cross-market spillover effect is an economic consequence of a severe power bottleneck. Investors allocate almost their entire inverter capacity to the highly profitable aFRR market, leaving only fractional MW capacities for day-ahead arbitrage. To still participate in the spot market, they massively oversize their cheap energy capacity (MWh), forcing the batteries into inefficient, slow-cycling, and sometimes loss-making spot market schedules. This fundamentally changes the traditional understanding of BESS revenue stacking, proving that optimal BESS capacities cannot be planned by evaluating single markets in isolation.

% Chapter \ref{ch:state_of_the_art} identified a distinct dichotomy in current literature, where studies typically focus strictly on either technical grid metrics or economic profit maximization. The novelty of this thesis was the development of a hybrid EPEC framework that co-optimizes multi-service market participation under strict physical grid constraints and battery degradation. The findings successfully confirm this promised progress beyond the state-of-the-art: By proving that technical constraints and economic multi-market bidding cannot be analyzed in isolation without distorting the investment case, this thesis provides a highly robust, integrated perspective that significantly advances current BESS allocation modeling.

% As critically discussed in Section \ref{sec:strengths_limitations}, the non-linear and non-convex nature of EPECs required pragmatic engineering compromises. Future work must aim to overcome the continuous approximations mandated by the current interior-point solver technology. The ultimate methodological milestone is the evolution from the current continuous formulation towards a MIPEC. Once commercial MI-NLP solvers achieve sufficient algorithmic maturity, future frameworks will be able to natively integrate discrete minimum bid sizes, true Pay-as-Bid auction clearing mechanisms, and complex non-differentiable battery degradation algorithms (e.g., Rainflow-Counting) without succumbing to mathematical singularities or locally infeasible points.

% Building upon the foundation laid in this thesis, several highly relevant avenues for future research emerge. First, to transition this framework into a comprehensive tool for macroeconomic forecasting, temporal upscaling is required. Expanding the model from a representative 24-hour period to an 8760-hour simulation -- or utilizing clustered seasonal representative days -- will allow for the evaluation of long-term seasonal phenomena, such as dunkelflauten.
% Second, the framework could be applied to investigate pressing European energy policy questions. By comparing the currently implemented \textit{Nodal Pricing} logic against European zonal pricing models, the model could quantify the economic inefficiencies of the current "copper plate" assumption and evaluate the much-debated bidding zone splits.
% Finally, future research should integrate \textit{co-located hybrid park configurations}. Allowing BESS units to share a grid connection point with wind or solar generators behind the meter would introduce a highly complex layer of strategy: co-optimizing grid-fee-free local charging against system-wide market participation. This would align the EPEC framework perfectly with the immediate realities of modern, large-scale renewable project development.

